\section*{Description of the Technology} \label{sec:tech_desc}
\noindent

[ACRONYM] is an approximate error correction scheme for data blocks stored in or transmitted through unreliable channels. Unlike algebraic codes (BCH, Reed-Solomon, LDPC), [ACRONYM] does not require finite-field arithmetic. Instead, it builds a set of CRC-based constraints over a two-dimensional arrangement of the data, and uses a graph-guided combinatorial search to find and correct bit-flip errors. The scheme is approximate in the sense that its correction guarantee is probabilistic rather than algebraic: for well-chosen parameters, the probability of full correction exceeds 95\% at 5\% BER, with the small failure probability attributable to hash collisions. In the next sections we describe the key components of [ACRONYM] in detail.

\section{Feistel Permutation}

\subsection{Purpose}

The first stage of [ACRONYM] is a keyed invertible permutation $\pi$ applied to the indices $\{0, 1, \ldots, L-1\}$ of the data block. This permutation serves two purposes: (i) it provides deterministic, reproducible reordering parameterized by a secret key $K$, shared between the encoder and decoder; and (ii) it scatters contiguous burst errors across the two-dimensional grid, converting the burst correction problem into the statistically equivalent random-error problem.

To see why (ii) holds, consider a burst of $b$ consecutive bit-flips at source positions $\{i, i+1, \ldots, i+b-1\}$. Under a uniformly random permutation, these positions are mapped to approximately uniformly distributed locations in $\{0, \ldots, L-1\}$, placing roughly one flipped bit per row and per column of the grid. The Feistel permutation approximates this uniform distribution while remaining invertible and keyed.

\subsection{Construction}

$\pi$ is implemented as an $R$-round balanced Feistel network, with $R = 8$ in the default configuration. Each round applies a keyed round function $F_K$ based on SHA-256: the input to SHA-256 is the concatenation of the right-half index value, the round index $r$, and the 16-byte key $K$; the output is truncated to the required bit width. The forward permutation is:

\begin{equation}
    \text{For } r = 0, \ldots, R-1: \quad R_i \leftarrow L_i \oplus F_K(R_i \,\|\, r), \quad L_i \leftarrow R_i
\end{equation}

where $L_i$ and $R_i$ are the left and right halves of the current index, and $\oplus$ denotes bitwise XOR.

Since $L$ is not generally a power of two, the Feistel network is extended to the exact domain $\{0, \ldots, L-1\}$ using the \emph{cycle-walking} technique: after the $R$ rounds, if the candidate output index $c$ falls outside $\{0, \ldots, L-1\}$, the cipher is re-applied to $c$, repeating until a valid index is obtained. The expected number of applications is at most 2 for typical $L$ values, and the resulting $\pi$ is provably bijective on $\{0, \ldots, L-1\}$.

The inverse permutation $\pi^{-1}$ is computed by applying the $R$ rounds in reverse order with swapped halves, followed by the same cycle-walking procedure. This allows the decoder to recover original bit positions from grid positions without storing any additional data.

\section{Two-Dimensional Grid Layout}

The permuted bit array $B' = \{B[\pi^{-1}(0)], B[\pi^{-1}(1)], \ldots\}$ is arranged into an $N \times N$ grid $G$, where $N = \lceil\sqrt{L}\rceil$. Grid cell $(r, c)$ holds bit $B'[r \cdot N + c]$ for $r \cdot N + c < L$, and zero for cells beyond position $L$ (padding). The grid is populated in row-major order.

For representative block sizes: $L = 256$ bits yields $N = 16$ (16$\times$16 grid); $L = 1{,}024$ bits yields $N = 32$; $L = 4{,}096$ bits yields $N = 64$; $L = 16{,}384$ bits yields $N = 128$.

\section{Hash Node Construction}

\subsection{Row Hash Nodes}

The $N$ rows of grid $G$ are divided into non-overlapping groups of $s$ consecutive rows ($s \geq 1$, default $s = 1$). For each group $g$, a row hash node is defined by: its group index, its axis (``row''), the set of source bit indices covered by the group (i.e., all source bits whose permuted position lies in those rows), and a CRC digest computed over the bits of those rows serialized in row-major order.

\subsection{Column Hash Nodes}

Analogously, the $N$ columns are divided into groups of $s$ consecutive columns. Each column group yields a column hash node with its CRC digest over the bits of those columns serialized in column-major order.

\subsection{Hash Functions}

[ACRONYM] supports a family of hash functions for computing digests:

\begin{itemize}
    \item \textbf{CRC-8} ($h=8$ bits): 25\% overhead at $L=4{,}096$; false positive probability $2^{-8} \approx 0.39\%$.
    \item \textbf{CRC-16} ($h=16$ bits): 50\% overhead; false positive probability $2^{-16} \approx 0.0015\%$.
    \item \textbf{CRC-32} ($h=32$ bits, polynomial \texttt{0xEDB88320}): 100\% overhead; false positive probability $\approx 2.3 \times 10^{-10}$.
    \item \textbf{CRC-64} ($h=64$ bits): 200\% overhead; false positive probability $\approx 5.4 \times 10^{-20}$.
    \item \textbf{GF(2) SimHash} ($h$ bits): configurable width; equivalent false positive probability to CRC of same width.
\end{itemize}

\subsection{Overhead Formula}

Total metadata is $2\lceil N/s \rceil \cdot h$ bits. The overhead ratio is:
\begin{equation}
    \text{overhead ratio} = \frac{2\lceil N/s \rceil \cdot h}{L} \approx \frac{2h}{s\sqrt{L}}
    \label{eq:overhead}
\end{equation}

This quantity decreases as $L$ increases for fixed $h$ and $s$, demonstrating the favorable $O(1/\sqrt{L})$ scaling of [ACRONYM].

\section{Constraint Graph}

After computing all hash nodes, [ACRONYM] constructs a constraint graph $G_C = (V, E)$ where $V$ is the set of all hash nodes and an edge $(u, v) \in E$ exists whenever $\text{covered\_indices}(u) \cap \text{covered\_indices}(v) \neq \emptyset$. Edge weight equals the intersection size.

In the default configuration ($s=1$), the constraint graph is a complete bipartite graph $K_{N,N}$ between the $N$ row nodes and $N$ column nodes. Each edge has weight 1, corresponding to the unique grid cell at the intersection of that row and column. This bipartite structure is the key to localization: a single bit-flip at grid position $(i,j)$ corrupts exactly one row node (row $i$) and one column node (column $j$), and the pair of mismatched nodes shares exactly one edge---pointing directly at the error location.

\section{Constraint-Guided Bit-Flip Search}

\subsection{Overview}

The decoder labels each hash node as matched (recomputed digest equals stored digest) or mismatched. It then iteratively selects mismatched nodes and searches for bit-flip combinations that increase the total number of matched nodes, using matched-neighbor pinning to prune the search.

\subsection{Matched-Neighbor Pinning}

For a mismatched node $u$ being processed, let $\mathcal{M}(u)$ denote the matched neighbors of $u$ and $\mathcal{X}(u)$ the mismatched neighbors. The decoder partitions the covered indices of $u$ into three disjoint sets:

\begin{itemize}
    \item \textbf{Intersection indices}: bits in $\text{covered}(u) \cap \bigcup_{v \in \mathcal{X}(u)} \text{covered}(v)$. These are implicated by multiple simultaneous violations.
    \item \textbf{Free indices}: bits in $\text{covered}(u)$ not covered by any matched neighbor.
    \item \textbf{Pinned indices}: bits in $\text{covered}(u) \cap \bigcup_{v \in \mathcal{M}(u)} \text{covered}(v)$. These are \emph{excluded} from the search.
\end{itemize}

The candidate set is $\mathcal{C}(u) = \text{intersection} \cup \text{free}$. Pinned indices are never flipped in the current iteration, preventing the decoder from corrupting already-correct constraints.

\subsection{Flip-Level Climbing}

For each mismatched node $u$, the decoder enumerates all combinations of $k$ indices from $\mathcal{C}(u)$ for $k = 1, 2, 3, \ldots$, up to a maximum $k_{\max}$. For each combination:
\begin{enumerate}
    \item Tentatively apply the $k$ bit-flips.
    \item Recompute hash digests for all affected nodes (at most $2k$ in the bipartite case).
    \item Score = total matched nodes in $V$.
    \item Revert and record the score.
\end{enumerate}

The combination achieving the highest score is committed. If no combination improves the score at the current $k$, $k_{\max}$ is incremented and the next pass begins. The outer loop repeats over all mismatched nodes until all are matched or no further progress is achievable.

\subsection{Decoding Output}

After correction, the corrected grid $G_c$ is serialized in row-major order and the inverse Feistel permutation $\pi^{-1}$ is applied to recover the corrected source bits in their original order.

\section{Empirical Performance}

Figure~\ref{fig:overhead} shows the overhead ratio versus block length for [ACRONYM] configurations (CRC-8, CRC-16, CRC-32) compared to BCH codes sized for the same correction capability at 5\% BER. [ACRONYM] with CRC-32 achieves 100\% overhead at $L=4{,}096$ versus BCH's $\sim$174\%, and overhead falls to $\sim$50\% at $L=16{,}384$ while BCH remains roughly constant.

Figure~\ref{fig:success} shows correction success rate versus flip count at $L=4{,}096$ bits for [ACRONYM] CRC-32, demonstrating $\geq$95\% full correction through 200 flips (5\% BER).

Figure~\ref{fig:burst} shows that [ACRONYM] achieves statistically identical correction success rates and search effort for random versus burst bit-flip patterns of equal total count, confirming the burst resilience provided by the Feistel permutation.

\begin{figure}[h]
    \centering
    % Place fig2 overhead comparison figure here
    % \includegraphics[width=0.85\textwidth]{fig/fig2_overhead.png}
    \caption{Overhead ratio vs. correctable errors: [ACRONYM] (CRC-8/16/32) compared to BCH analytical bounds at $L=4{,}096$ bits. [ACRONYM] CRC-32 achieves 100\% overhead with $\geq$200 correctable flips; BCH requires $\sim$174\% for equivalent capability.}
    \label{fig:overhead}
\end{figure}

\begin{figure}[h]
    \centering
    % Place fig1 success rate figure here
    % \includegraphics[width=0.85\textwidth]{fig/fig1_success.png}
    \caption{Correction success rate vs. flip count at $L=4{,}096$ bits for [ACRONYM] CRC-32 (100\% overhead). Full correction is achieved for 200+ flips ($\geq$5\% BER) with $\geq$95\% success rate.}
    \label{fig:success}
\end{figure}

\begin{figure}[h]
    \centering
    % Place fig4 burst resilience figure here
    % \includegraphics[width=0.85\textwidth]{fig/fig4_burst.png}
    \caption{Correction success rate (left) and solver search effort (right) for random versus burst bit-flip errors at $L=4{,}096$ bits. Curves overlap, confirming that the Feistel permutation renders [ACRONYM] burst-resilient with no additional overhead.}
    \label{fig:burst}
\end{figure}
